\scanpage{32}
rang $1$.
\[
\Sigma^2(F)/\mathbb V(L^{\otimes2})
\]
est canoniquement isomorphe à $2F\subset\check{\mathbb V}\check B=\mathbb V(B)$
image inverse dans $\check{\mathbb V}(\check B)$ de la section $2$
de $\check{\mathbb V}(\underline O_S)$, i.e
\[
\Sigma^2(F)/\mathbb V(L^{\otimes2})\simeq 2F \tag*{(39)}
\]

\deleted{Ainsi} La donnée d'une section $\sigma_{\unsure{*|\bullet}}$ de
$\Sigma^2(\deleted{\unsure{\check{\mathbb V}}}F)$, ou encore d'un diviseur quadratique
$\Omega\subset\deleted{\unsure{\check B}}F$, \deleted{\unsure{q}} définit donc une section de
$2F$, qu'on \deleted{appellera la fonction} dénote par
$b_\Omega$. \deleted{Et} Le morphisme
\[
\Sigma^2(F)\xrightarrow{\underline b}2F \tag*{(40)}
\]
est caractérisé par la condition de commuter
aux chgts de base, et de transformer
la section de $\Sigma^2(F)$ image d'une section
$(s_1,s_2)$ de $F^2$, en la section $s_1+s_2$ de $2F$
(ainsi $\underline b$ est une sorte de \uline{trace universelle});
en \unsure{terme|termes} de $\Omega$, et on peut aussi dire
que $\underline b$ (interprété en termes de
$\Omega\longmapsto \underline b_\Omega$) est défini par la condition de
compatib. aux chgts de base, \uline{et} \deleted{que} de
\deleted{transformer} le \deleted{\unsure{deux sections}} diviseur quadratique

% archive sheet number: 31

